\section{评价指标与实验协议}
\subsection{评价对象、参考速度修正与记录量}
同时保存原始命令、MLP 均值、采样修正、处理后修正、参考速度修正、实际机身速度、参考模型状态、逐子步受力、接触、姿态和终止原因。以实际冲量除以控制周期构造 $\bar{\bm w}_k$，按相同时间基准计算所有算法的动态目标。

模型一致性指标使用预先固定的 $\bm M,\bm D$ 及外力序列积分得到独立的参考速度修正。参考模型从 episode 的统一初态初始化，避免每步用实际状态重置而掩盖长期偏差。全身实验另外记录末端速度、位置、腕力和工作空间约束。算法内部 reward 作为训练诊断保留。

\subsection{拟合、运动与动力学一致性}
\begin{longtable}{P{28mm}P{84mm}P{30mm}}
\toprule 指标 & 定义与计算方式 & 单位/解释\\\midrule\endhead
修正 RMSE & 各轴 $\sqrt{N^{-1}\sum_k(\delta_{j,k}-\delta_{j,\rm ref,k})^2}$，分别报告均值输出和处理后输出 & m/s 或 rad/s\\
实际速度 RMSE & 各轴 $\sqrt{N^{-1}\sum_k(v_{j,k}-c^0_{j,k}-\delta_{j,\rm ref,k})^2}$，转向同理 & 反映闭环响应\\
稳态力--速度增益 & 在无饱和恒力窗口拟合 $\bar\delta_j=\hat g_jF_j+b_j$，比较 $\hat g_j$ 与 $1/d_j$ & m/(N\,s)，报告偏置及置信区间\\
响应时间常数 & 恒力或释放段拟合一阶模型，比较 $\hat T_j$ 与 $m_j/d_j$ & s，惯量响应\\
动力学残差 & 单轴 $e_{F,k}=m(\delta^{\rm act}_{k+1}-\delta^{\rm act}_k)/\Delta t+d\delta^{\rm act}_k-\bar F_k$ & N，偏航为 N\,m\\
端点/全身误差 & 末端位置 RMSE、速度 RMSE、腕力 RMSE，按固定坐标系分轴报告 & m、m/s、N\\\bottomrule
\caption{第一组指标。$\delta^{\rm act}=v_{\rm actual}-c^0$，辨识时保持原始命令恒定。}
\end{longtable}

动力学残差采用所有方法一致的滤波与差分规则，同时披露滤波带宽和延迟。力--速度增益与时间常数在未饱和区间辨识，饱和区间单独报告。对于当前三值参考速度修正，在非零区间内，其对力的局部斜率为零，这可以由公式证明，实际机器人斜率仍需测量。

\subsection{牵引起停、交互与稳定性}
\begin{longtable}{P{28mm}P{84mm}P{30mm}}
\toprule 指标 & 定义与计算方式 & 单位/解释\\\midrule\endhead
牵引响应延迟 & 从真实受力开始，到沿预期方向的实际速度修正首次超过预设阈值并持续 0.2 s & s，阈值随协议公布\\
松手恢复时间 & 从实际受力消失，到 $\|\delta^{\rm act}_{xy}\|<0.05$ 且 $|\delta^{\rm act}_{\omega}|<0.05$，持续 0.5 s & s，两阈值单位分别为 m/s、rad/s\\
松手残余运动 & 固定释放后窗口内积分 $\int\|\delta^{\rm act}_{xy}\|dt$ 与 $\int|\delta^{\rm act}_{\omega}|dt$，同时比较目标衰减轨迹 & m、rad，目标依参数确定\\
无力误响应 & 无力片段中，处理后修正超过上述阈值的时间占比与 RMS & \%，m/s、rad/s\\
过冲与平滑度 & 目标非零时的归一化峰值过冲，实际加速度/jerk 的 RMS 与 P95，修正变化率 & \%、m/s$^2$、m/s$^3$\\
碰撞峰值与冲量 & 固定障碍测试的 $\max\|F_{\rm contact}\|$、$\int\|F_{\rm contact}\|dt$，统一碰撞前速度 & N、N\,s\\
稳定生存率 & 完成固定时长且无环境定义的跌倒、禁触和约束终止的 episode 比例，所有方法使用同一终止配置 & \%，Wilson 95\% CI\\
姿态与执行约束 & roll/pitch 的峰值与 P95，关节越界占比、力矩饱和占比、足滑速度 & rad、\%、m/s\\
计算与能耗 & 目标硬件上端到端推理 P50/P95，平均 $\sum_i|\tau_i\dot q_i|$ & ms、W\\\bottomrule
\caption{第二组指标。提前终止或未恢复的 episode 计为失败/超时，同时给出数量，避免只统计成功片段。}
\end{longtable}

任务成功率另外要求 episode 存活、实际速度对由原始命令与独立参考速度修正叠加得到的目标速度的平移 RMSE 小于 $0.15\,\mathrm{m/s}$、转向 RMSE 小于 $0.15\,\mathrm{rad/s}$，且释放恢复时间超过目标模型恢复时间的部分不大于 $0.5\,\mathrm s$。这些容差作为初版实验设计预先固定，并附容差敏感性曲线。

具有虚拟惯量时，释放后的非零速度属于目标动力学的一部分。恢复时间应与给定 $m/d$ 的目标曲线一起评价，额外报告相对目标的释放段积分误差。正向持续命令下，以相对原始命令的速度修正衡量恢复，避免把正常继续行走计为漂移。静止场景直接报告平均绝对功率，避免用接近零的位移计算能耗距离比。

\subsection{测试矩阵与消融}
以下为建议采用的实验协议，数值属于待执行设计。每个 episode 长 20 s，单次牵引从第 2 s 开始，按测试条件持续 0.2、1 或 3 s，剩余时间观测释放恢复。每个随机方法至少训练 5 个独立种子，每个测试条件至少评估 100 个 episode，使用一致的底座 checkpoint、环境种子和牵引时间序列。报告种子级均值与 95\% bootstrap 区间，并按场景展示 P95。训练集与测试集按完整轨迹、地形实例和扰动种子拆分。

\begin{table}[htbp]\centering\small
\begin{tabularx}{\linewidth}{P{29mm}Y}
\toprule 维度 & 覆盖条件\\\midrule
原始命令 & 零命令、直行、侧行、原地转向，包含训练范围内与保留的测试组合\\
力幅值与方向 & $0,5,10,15,17,17.5,18,20,30,40,50\,\mathrm N$，正反方向与 45 度方向，补充负 $x$ 轴两侧方向\\
时间波形 & 持续 0.2/1/3 s 的阶跃，斜坡、频率扫描、阈值附近波动、突然释放与反向牵引\\
地形 & 平地、坡地、粗糙地形、上下台阶，记录实际坡度、台阶高和粗糙幅值，同时保留地形等级\\
泛化 & 质量/摩擦、传感噪声、动作延迟和牵引作用点变化，记录参数数值与分布\\
可调参数 & 以一个经验证可执行的基准 $(m_0,d_0)$ 为中心，分别测试 $0.5/1/2$ 倍，含训练内与留出组合\\\bottomrule
\end{tabularx}
\caption{评价测试矩阵。高幅值边界测试按机器人仿真与硬件能力确定执行范围。}
\end{table}

地形等级是课程变量，不能独立表征实际地形难度。论文中的比较应固定地形生成种子和几何参数，同时报告等级分布。大角度牵引及大力实验优先验证坐标变换、动作限幅和 episode 终止条件的一致性。

\begin{table}[htbp]\centering\small
\begin{tabularx}{\linewidth}{P{13mm}Y Y}
\toprule 组别 & 配置 & 要回答的问题\\\midrule
A & 原始 83 维 + 当前参考速度修正 + PPO & 当前配方的测量基线\\
B & 80 维去冗余 + 当前参考速度修正 + PPO & 三维冗余的影响\\
C & 原始 83 维 + 当前参考速度修正 + MSE & 标签拟合训练是否足够\\
D & 80 维 + 动态参考速度修正 + MSE & 动态参考速度修正能否保留连续幅值与时序\\
E & 与 D 相同输入及参考速度修正 + PPO 或 MSE 后 PPO & 闭环强化学习的增量收益\\
F & 真值力 + 显式导纳 + 同一冻结底座 & 底座可实现响应与估计误差的分解\\
G & 监督/GRU 估计 + 显式导纳 + 同一底座 & 无力传感条件下的可部署性能\\\bottomrule
\end{tabularx}
\caption{逐因素消融。B 对 A 检验输入，C 对 A 检验算法。参考速度修正的设计另需在相同 80 维和训练算法下补齐静态/动态成对实验。}
\end{table}

训练预算报告环境步数、梯度更新次数和墙钟时间。F 使用特权力，作为参考性能条件单列。MSE 组应覆盖自身闭环访问的数据分布。各组统一报告动态一致性、无力误响应和稳定性指标，再分别报告各自参考速度修正的拟合误差，避免定义变化影响指标可比性。
